% !PRINTED 390
\weberpage{390}
Let us now consider any primitive quadratic equation with
integer coefficients
\weq{16}{A + B\Omega + C\Omega^2 = 0,}
in which $A$, $B$, $C$ are integers without common divisor,
satisfying the condition
\[
B^2 - 4AC = D
\]
The two roots $\Omega$, $\Omega'$ of this equation are quadratic
irrational numbers belonging to the discriminant $D$, and
if we develop them in continued fractions, then by
\S.\,124 we shall at last come upon closing numbers $\omega_1$, $\omega_2$ which belong to the reduced
numbers, and which must therefore be contained in the periods
discussed above.

It remains to determine whether these closing numbers $\omega_1$, $\omega_2$ are
contained in the same period or in different periods.

By 6. they are contained in the same period if and only
if $\Omega$ and $\Omega'$ are equivalent to each other. If we call
such quadratic irrational numbers as are equivalent to their conjugates
\gesperrt{two-sided numbers}\footnote{``\gesperrt{Two-sided}'' we use, following \gesperrt{Dedekind}'s proposal, in place of the
\gesperrt{Gauss}ian \gesperrt{anceps} or the otherwise usual \gesperrt{ambig}: \gesperrt{Dirichlet-Dedekind}, Vorlesungen über Zahlentheorie. 4th edition, 1894, p. 139.}, then we can
say that the continued-fraction development of the roots of
equation (16) leads to one or to two different periods,
according as the roots are two-sided or not.
Which of the two cases occurs can be recognised from the period
itself.

If, namely, $\omega$ is a reduced number, and $\omega'$ conjugate
to $\omega$, then $-1 : \omega'$ is also reduced; and if $\Omega$ is equivalent
to $\omega$, then $\Omega'$ is equivalent to $\omega'$, and hence also to $-1 : \omega'$ (\S.\,120).

If
\[
\omega = \{a, b, c\},
\]
then by (12):
\[
-\frac{1}{\omega'} = \{c, b, a\}.
\]

By \S.\,124, (9), $-1 : \omega'$ therefore follows $-1 : \omega_1'$
in the same way as $\omega_1$ follows $\omega$. If therefore the period of $\omega$ is
\[
\omega, \omega_1, \omega_2 \ldots \omega_{\nu-1},
\]
then the period of $-1 : \omega'$ is:
\[
\frac{-1}{\omega'_{\nu-1}}, \frac{-1}{\omega'_{\nu-2}} \cdots \frac{-1}{\omega_1'}, \frac{-1}{\omega'}.
\]
